\chapter{Einleitung}

\section{Motivation und Zielsetzung}
%faik
Die Entwicklung mechatronischer Systeme bietet Studierenden die Möglichkeit, theoretisch erworbenes Wissen praxisnah anzuwenden und interdisziplinäre Zusammenhänge zwischen Mechanik, Elektronik und Informatik zu verstehen. Im Rahmen dieser Studienarbeit soll ein mechatronischer Demonstrationsroboter entstehen, der auf Hochschulmessen eingesetzt wird, um das Studienangebot der DHBW und die Vielfältigkeit ingenieurwissenschaftlicher Tätigkeiten anschaulich zu präsentieren. Der Roboter soll durch seine technische Funktionsvielfalt und sein interaktives Auftreten insbesondere das Interesse technisch affiner Besucherinnen und Besucher wecken und die Attraktivität des Studiengangs Fahrzeugelektronik hervorheben.

Das Ziel der Arbeit besteht in der Konzeption, Entwicklung und Realisierung eines modular aufgebauten Demonstrators, der grundlegende Bewegungsfunktionen ausführt und über eine Fernsteuerung oder eine mobile Applikation bedient werden kann. Ergänzend soll der Roboter über ein Display verfügen, das Informationen zur DHBW und zum Studiengang anzeigt, sowie über akustische Ausgabemöglichkeiten zur Interaktion mit dem Publikum. Das Design und die Systemarchitektur werden so ausgelegt, dass eine zukünftige Erweiterung um autonome Funktionen, etwa eine sensorbasierte Hinderniserkennung oder KI-gestützte Bewegungssteuerung, problemlos möglich ist. Auf diese Weise entsteht ein nachhaltiger Demonstrator, der nicht nur für aktuelle Messeauftritte genutzt, sondern auch als Basis für weiterführende studentische Projekte dienen kann.

\section{Aufgabenstellung}
%faik
Die Aufgabenstellung umfasst die Planung und Umsetzung aller erforderlichen Schritte zur Realisierung des Demonstrators. Dazu gehören die Analyse vergleichbarer bestehender Systeme, die Definition technischer Anforderungen sowie die Ausarbeitung eines modularen Systemkonzepts. Anschließend erfolgt die Auswahl und Integration der mechanischen, elektrischen und softwareseitigen Komponenten.
Ein Schwerpunkt liegt auf der Entwicklung einer zuverlässigen Antriebs- und Steuerungsarchitektur sowie einer intuitiven Benutzeroberfläche für die Fernbedienung bzw. App-Steuerung. 
Nach der mechanischen und elektronischen Umsetzung folgen Inbetriebnahme, Funktionsprüfung und Dokumentation der erzielten Ergebnisse.  
Ziel ist es, einen funktionsfähigen, erweiterbaren Roboter bereitzustellen, der für den Messeeinsatz geeignet ist und zukünftige Weiterentwicklungen ermöglicht.


\section{Anforderungen an den Demonstrator(lieber weglassen?)}
%faik
Für den geplanten Messeeinsatz ergeben sich sowohl technische als auch funktionale Anforderungen. Der Demonstrator soll ein kompaktes, robustes und transportfreundliches Design aufweisen und eine sichere, energieeffiziente Betriebsweise ermöglichen. 
Funktional sind grundlegende Bewegungsabläufe wie Vorwärts-, Rückwärts- und Drehbewegungen zu realisieren, gesteuert über eine Fernbedienung oder mobile Anwendung. 
Zur Interaktion mit dem Publikum sind ein Display zur Informationsanzeige sowie akustische Ausgabemöglichkeiten vorgesehen. 
Darüber hinaus ist die Systemarchitektur so zu gestalten, dass spätere Erweiterungen - etwa autonome Bewegungsfunktionen, zusätzliche Sensorik oder KI-basierte Steuerungsmechanismen - problemlos integriert werden können. 
Neben diesen technischen Zielen wird auch auf Wartbarkeit, Modularität und Kosteneffizienz geachtet\footnote{Für eine detaillierte Anforderungsliste siehe Abschnitt~\ref{sec:Projekt Requierments}.}.

\section{Aufbau der Arbeit}
%faik
Zu Beginn der Arbeit erfolgt die Entwicklung eines grundlegenden Konzeptes für den Demonstrationsroboter. Dieses entsteht im Rahmen gemeinsamer Brainstorming-Sitzungen und Abstimmungen innerhalb des Projektteams. Dabei werden erste Funktionsideen, Designvorschläge und Einsatzszenarien diskutiert und bewertet. Anschließend folgt eine systematische Recherche, um die technische und wirtschaftliche Umsetzbarkeit des Konzepts zu überprüfen. Hierbei werden bestehende Lösungen analysiert, potenzielle Schwierigkeiten identifiziert und Optimierungsmöglichkeiten hinsichtlich Funktionalität, Kosten und Realisierbarkeit untersucht.

Auf Grundlage dieser Erkenntnisse wird das ursprüngliche Konzept schrittweise konkretisiert und an die verfügbaren Ressourcen, Kompetenzen und wirtschaftlichen Rahmenbedingungen angepasst. Danach erfolgt die Auswahl, Beschaffung und Erprobung der erforderlichen Komponenten. Die Entwicklung der einzelnen Funktionsmodule wird zunächst separat durchgeführt, um eine gezielte Fehleranalyse und Funktionsprüfung zu ermöglichen. Im Anschluss werden alle Teilfunktionen in das Gesamtsystem integriert und getestet, sodass am Ende ein vollständig funktionsfähiger Demonstrator entsteht.




%\label{cha:Einleitung}
%
%Folgende Stichworte können zum Aufbau der Einleitung herangezogen werden.
%
%\begin{itemize}
%\item Hinführung, Begründung, Zweck und Ziel der Aufgabenstellung
%\item Erläuterung der Problemstellung
%\item Konkretisierung der zu lösenden Aufgabe
%\item Gegebenenfalls Formulierung einer Leitfrage oder Forschungsfrage
%\item Ausgangslage, geplante Vorgehensweise, Methoden zur Bearbeitung und Zielsituation
%\item Zum Ende der Einleitung wird eine Kurzübersicht über die Inhalte der Kapitel gegeben: \glqq Die Arbeit ist wie folgt gegliedert: ...\grqq
%\end{itemize}
%
%Die Einleitung wird üblicherweise auf ein bis zwei Seiten als fortlaufender Text geschrieben. Eine weitere Untergliederung in nummerierte Abschnitte ist nicht empfehlenswert, da dies erstens unüblich ist, zweitens die Lesbarkeit nicht begünstigt und drittens die Formulierung der Einleitung erschwert. Weitere Empfehlungen zum Aufbau der Einleitung und des gesamten Dokuments sind z.~B. aus \autocite{DHBW.2021} und \autocite{Lindenlauf.2022} zu entnehmen.
%
%\clearpage
%
%Hinweise: 
%
%\begin{itemize}
%\item Auch in der Einleitung unbedingt zu wichtigen Hintergründen und Fakten Zitate aufführen.\index{Zitat} Zitate bitte in der Form \autocite{Tipler.2019} oder mit Seitenbezug \autocite[66]{Ziegler.2017} oder auch mehrere Zitate  \autocite{Tipler.2019, Ziegler.2017} innerhalb einer eckigen Klammer angeben. Zur besseren Lesbarkeit bitte immer ein Leerzeichen vor dem Zitat einfügen.
%	
%\item Bereits in der Einleitung können Abkürzungen erläutert werden. Grundsätzlich gilt, dass bei der ersten Verwendung einer Abkürzung diese auch erläutert wird. Zum Beispiel können das \acrlong{acr:ABS}  (\acrshort{acr:ABS}) oder die Fahrdynamikregelung (\acrlong{acr:ESC}, \acrshort{acr:ESC}) als Abkürzungen eingeführt werden. In der Datei \textit{pages/abkuerzungen.tex} sind alle verwendeten Abkürzungen einzufügen. Neben dem verpflichtenden Abkürzungsverzeichnis kann auch ein Glossar hinzugefügt werden. In dieser Vorlage können Glossareinträge in der Datei \textit{pages/glossar.tex} eingefügt werden. Ein \Gls{gls:glossar} ist jedoch nicht verpflichtend.
%\end{itemize}